\subsection{Sortieralgorithmen}

\subsubsection{SelectionSort}
Sucht das kleinste Element in der Liste, fügt es am Anfang einer neuen Liste ein und löscht es aus der alten Liste. Wiederholt, bis alte Liste leer ist $\rightarrow O(n^2)$

\subsubsection{InsertionSort}
\begin{enumerate}
    \item Wähle num = input[i]
    \item Vergleiche alle Indices j $<$ i und schiebe num nach input[j], wenn input[j] $<$ num
\end{enumerate}
\begin{tabularx}{\linewidth}{@{}>{\bfseries}l@{\hspace{.5em}}|X@{}}
    Laufzeit & \textbf{Best case}: $O(n)$ falls Input sortiert \\
             & \textbf{Worst case}: $O(n^2)$ falls Input absteigend sortiert \\
             & $O(hn)$ falls ein Element höchstens $h$ Elemente von seiner Zielposition entfernt ist
\end{tabularx}

\subsubsection{MergeSort}
\textbf{Divide and Conquer - Rekursiv}
\begin{enumerate}
    \item Liste wird rekursiv halbiert, bis zwei Listen mit je einem Element entstehen
    \item Die beiden Elemente werden verglichen und sortiert in die Rückgabeliste eingefügt: \@returns Liste mit 2 sortierten Elementen
    \item Zwei Listen mit 2 Elementen werden sortiert gemerged, indem list1[i] mit list2[i] verglichen wird und das kleinere der beiden in die Ergebnisliste geschrieben wird
    \item Wiederholung für alle Teillisten, bis am Ende zwei Listen mit $n/2$ sortierten Elementen gemerged werden $\rightarrow$ sortierte Liste
\end{enumerate}
\begin{tabularx}{\linewidth}{@{}>{\bfseries}l@{\hspace{.5em}}|X@{}}
    Laufzeit & \textbf{Worst case}: $O(n\log_2 n)$
\end{tabularx}

\subsubsection{QuickSort}
\textbf{Divide and Conquer - Rekursiv}
\begin{enumerate}
    \item Wähle ein Pivotelement A[p]
    \item Bewege alle Elemente < A[p] nach links von A[p]
    \item Bewege alle Elemente > A[p] nach rechts von A[p]
\end{enumerate}
\begin{tabularx}{\linewidth}{@{}>{\bfseries}l@{\hspace{.5em}}|X@{}}
    Laufzeit & \textbf{Worst case}: $O(n^2)$ falls als Pivotelement immer das größte Element gewählt wird (sehr selten)
\end{tabularx}

\subsection{Suchalgorithmen}
\subsubsection{Binary Search}
\textbf{Annahme}: Liste sortiert, gesucht $x$.
\begin{enumerate}
    \item Vergleiche $x$ mit dem Wert, der in der Mitte steht
    \begin{itemize}
        \item[a)] Bei Gleichheit sind wir fertig
        \item[b)] Falls $x$ kleiner als der mittlere Wert, gehe zu $1.$ mit dem linken Teil des Arrays
        \item[c)] Falls $x$ größer als der mittlere Wert, gehe zu $1.$ mit dem rechten Teil des Arrays
    \end{itemize}
\end{enumerate}

\begin{tabularx}{\linewidth}{@{}>{\bfseries}l@{\hspace{.5em}}|X@{}}
    Laufzeit & \textbf{Best case}: $O(1)$ falls mittlere Wert gleich gesuchter Wert \\
             & \textbf{Worst case}: $O(\log_2 n)$
\end{tabularx}